% Part: sets-functions-relations
% Chapter: arithmetization
% Section: cauchy

\documentclass[../../../include/open-logic-section]{subfiles}

\begin{document}

\olfileid{sfr}{arith}{cauchy}

\olsection{ضميمه: حقيقي عددونه د کوشي د لړيو په توګه}

په \olref[cuts]{sec} کښې مو حقيقي عددونه د ډېډېکېنډ د پرېکونونو په
توګه جوړ کړل۔ په دې برخه کښې يوه بله جوړونه تشريح کوو۔ دا د کوشي پر
هغه تعريف ولاړه ده چې اوس ئې د کوشي لړۍ بولو؛ خو له دې تعريف څخه د
حقيقي عددونو د \emph{جوړولو} کار د نولسمې پېړۍ نورو ليکوالانو، په
ځانګړې توګه وايرشټراس، هاينې، مېراے او کانټور، کړے دے۔ (د ښۀ تاريخ
دپاره \citeauthor{OConnorRobertson:RN} \citeyear{OConnorRobertson:RN}
وګورئ۔)

مخکښې له دې چې نولسمې پېړۍ ته ورسېږو، د سايمن سټېوين (1548--1620)
خبره د پام وړ ده۔ په لنډ ډول، سټېوين پوه شو چې هر حقيقي عدد د هغه د
اعشاري پراختيا له مخې ګڼلے شو۔ نو آن يو غير ناطق عدد لکه $\sqrt{2}$
هم يوه ښۀ اعشاري پراختيا لري، چې پيل ئې داسې دے:
\[
	1.41421356237\ldots
\]
په سټ نظريه کښې د اعشاري پراختياوو نمونه جوړول ډېر اسان دي: يوازې
داسې تابعې $d \colon \Nat \to \Nat$ ئې وګڼئ چې $d(n)$ هغه $n$م
اعشاري ځاے وي چې موږ ئې څېړو۔ بيا لږ بدلون ته اړتيا لرو، څو د حقيقي
عدد هغه برخه هم سمبال کړو چې له اعشاري ټکي مخکښې راځي (دلته يوازې
$1$)۔ يو بل بدلون، يعنې د هم ارزښتۍ اړيکه، هم پکار ده څو د بېلګې په
توګه يقيني کړو چې $0.999\ldots = 1$۔ خو په سټ نظريه کښې د سټېوين په
لار د حقيقي عددونو بشپړه دقيقه جوړونه وړاندې کول ګران نۀ دي۔

سټېوين زموږ اصلي موضوع نۀ ده۔ (د سټېوين په اړه نور معلومات په
\citealt{KatzKatz2012} کښې وګورئ۔) خو دا ډېره ورته خبره وګورئ۔ د
$\sqrt{2}$ اعشاري پراختيا په نېغه د څېړلو پر ځاے، د $\sqrt{2}$ د
لا کره پراختياوو په اخيستلو د هغو ناطقو اټکلونو يوه \emph{لړۍ}
کتلے شو چې کره‌والے ئې پرله‌پسې زياتېږي:
\[
1, 1.4, 1.414, 1.4142, 1.41421,\ldots
\]
دا نظر چې حقيقي عددونه د «پرله‌پسې ښو اټکلونو» له لارې ګڼلے شو، د
يوې بلې پوهېدنې بنسټ راکوي (هغو پوهېدنو ته ورته چې $\Rat$ مو له
$\Int$، يا $\Int$ مو له $\Nat$ څخه جوړول)۔ بنسټيزې خبرې دا دي:
\begin{enumerate}
	\item هر حقيقي عدد د يوې (کېدے شي نامتناهي) اعشاري پراختيا په
	بڼه ليکل کېدے شي۔
	\item په يوه (کېدے شي نامتناهي) اعشاري پراختيا کښې زېرمه شوي
	معلومات په هماغه ډول د ناطقو عددونو په يوه لړۍ کښې زېرمه کېدے شي۔
	\item د ناطقو عددونو لړۍ له $\Nat$ څخه $\Rat$ ته يوه تابع ګڼلے
	شو؛ يوازې $f(n)$ په لړۍ کښې $n$م ناطق عدد وټاکئ۔
\end{enumerate}
البته له $\Nat$ څخه $\Rat$ ته هره \emph{تابع} حقيقي عدد نۀ راکوي۔
د بېلګې په توګه دا تابع وګورئ:
\[
	f(n) =\begin{cases}
			1 & \text{که }n\text{ طاق وي}\\
			0 &\text{که }n\text{ جفت وي}
		\end{cases}
\]
اصلي اندېښنه دا ده چې د $0,1,0,1,0,1,0,\ldots$ لړۍ هېڅ حقيقي
عدد ته «نژدې کېدونکې» نۀ ښکاري۔ نو د دې د يقيني کولو دپاره چې هغه
لړۍ څېړو چې کوم حقيقي عدد ته ورنژدې کېږي، بايد خپل پام هغو لړيو ته
محدود کړو چې کوم \emph{حد} لري۔

موږ د حد له نظر سره په \olref[his][set][limits]{sec} کښې مخکښې هم
مخامخ شوي يو۔ خو کټ مټ هماغه پخوانے تعريف نشو کارولے۔ هلته عبارت
«$(\forall \epsilon>0)$» په پټه پر حقيقي عددونو کميت‌ايښودنه
درلوده؛ او موږ د حقيقي عددونو پر تابعو حدونه څېړل؛ نو د هغه تعريف
کارول به دا وي چې حقيقي عددونه لا له مخکښې ومنو، حال دا چې موږ کټ
مټ د هغو \emph{جوړول} غواړو۔ نېکمرغي ده چې د حد له يوه ډېر نژدې
نظر سره کار کولے شو۔
\begin{defn}\ollabel{def:CauchySequence}
  يوه تابع $f: \Nat \to \Rat$ هله او يوازې هله د \emph{کوشي لړۍ}
  ده چې د هر مثبت $\epsilon \in \Rat$ دپاره
  $(\exists \ell \in \Nat)(\forall m, n > \ell)|f(m) - f(n)| <
  \epsilon$ وي۔
\end{defn}

د حد عمومي نظر د پخوا په شان دے: که د کره‌والي يوه ټاکلې کچه
(چې په~$\epsilon$ اندازه کېږي) وغواړئ، د کتلو يوه «سيمه» شته
(هره ورودي چې له~$\ell$ څخه لويه وي)۔ دا هم په اسانۍ ليدل کېږي چې
زموږ لړۍ $1$، $1.4$، $1.414$، $1.4142$، $1.41421$\ldots حد لري:
که غواړئ $\sqrt{2}$ د $\nicefrac{1}{10^{n}}$ له خطا دننه اټکل
کړئ، نو يوازې له $n$م غړي وروسته هر غړے وګورئ۔

نو څرګند فکر به دا وي چې حقيقي عدد هماغه هره د کوشي لړۍ \emph{ده}۔
خو لکه د $\Int$ او $\Rat$ په جوړونو کښې، دا خبره به بېخي ساده وي:
د هر ورکړل شوي حقيقي عدد دپاره څو بېلې د کوشي لړۍ هماغه حقيقي عدد
ښيي۔ د دې د ليدلو يوه اسانه لار داسې ده۔ يوه د کوشي لړۍ~$f$ واخلئ
او $g$ کټ مټ د~$f$ په شان تابع تعريف کړئ، يوازې دومره توپير چې
$g(0)\neq f(0)$۔ ځکه دواړه لړۍ له لومړي عدد وروسته په هر ځاے کښې
سره برابرې دي، نو په پاے کښې به ووايو چې د
\olref{def:CauchySequence} په معنا يو حد لري، او بايد د هماغه حقيقي
عدد «تعريفوونکې» وګڼل شي۔ نو په رښتيا بايد دا د کوشي لړۍ هماغه يو
حقيقي عدد وګڼو۔

په پايله کښې بيا بايد د کوشي پر لړيو د هم ارزښتۍ يوه اړيکه تعريف
کړو او حقيقي عددونه د هم ارزښتۍ له ټولګيو سره يو وګڼو۔
\textbf{د ژباړې د نسخې سمون (OLARI-010):} انګرېزي سرچينه په دې
ځاے «د هم ارزښتۍ اړيکې» وايي، خو جوړېدونکي حقيقي عددونه د اړيکې
\emph{ټولګي} دي؛ د لاندې تعريف او فورمول سره سم دلته هماغه ټولګي
ليکل شوي۔ لومړے داسې تابع پکار ده چې په حد کښې $0$ ته ورنژدې شي۔
د هرې تابع $h : \Nat \to \Rat$ دپاره ووايئ چې \emph{$h$ $0$ ته
ورنژدې کېږي} هله او يوازې هله چې د هر مثبت $\epsilon \in \Rat$
دپاره $(\exists \ell \in \Nat)(\forall n > \ell)|h(n)| <
\epsilon$ وي۔\footnote{دا د $\lim_{x
\mathord{\rightarrow}\infty}f(x) = 0$ له تعريف سره پرتله کړئ چې په
\olref[his][set][limits]{sec} کښې راغلے۔} برسېره پر دې، چې $f$ او
$g$ د $\Nat \to \Rat$ تابعې وي، $(f-g)(n) = f(n) - g(n)$
وټاکئ۔ اوس تعريف کړئ:
\[
	f \Realequiv g \text{ هله او يوازې هله چې $(f-g)$ $0$ ته ورنژدې کېږي}.
\]
بايد وارزوو چې $\Realequiv$ د هم ارزښتۍ اړيکه ده؛ او رښتيا هم ده۔
بيا، که وغواړو، حقيقي عددونه د $\Nat \to \Rat$ د ټولو کوشي
لړيو هغه د هم ارزښتۍ ټولګي تعريفولے شو چې د $\Realequiv$ له مخې
جوړېږي۔

\begin{prob}
د هر $n$ دپاره $f(n) = 0$ وټاکئ۔ او $g(n) = \frac{1}{(n+1)^2}$
وټاکئ۔ وښيئ چې دواړه د کوشي لړۍ دي او په رښتيا د دواړو تابعو حد
$0$ دے، نو $f \sim_\Real g$ هم دے۔
\textbf{د ژباړې د نسخې سمون (OLARI-011):} دلته سرچينه د تازه
تعريف شوې نښې پر ځاے بېله نښه کاروي؛ د تمرين معنا د هماغې تازه تعريف شوې اړيکې له
مخې هم ارزښت دے۔
\end{prob}

له دې وروسته، د معمول په شان، $\equivrep{f}{\Realequiv}$ د هغه د
هم ارزښتۍ ټولګي دپاره ليکو چې $f$ ئې يو !!a{element} دے۔ خو د لوست
د اسانولو دپاره به په راتلونکو برخو کښې لاندې‌ليک پرېږدو او يوازې
$\equivrep{f}{}$ ليکو۔ دا هم ټاکو چې د هر $q \in \Rat$ دپاره
$q_{\Real} = \equivrep{c_{q}}{}$ وي، چې $c_{q}$ ثابته تابع ده او
د هر $n \in \Nat$ دپاره $c_q(n) = q$ وي۔ بيا پر حقيقي عددونو
بنسټيزې اړيکې او عمليې داسې تعريفوو، د بېلګې په توګه:
\begin{align*}
	\equivrep{f}{} + \equivrep{g}{} &= 	\equivrep{(f + g)}{} \\
	\equivrep{f}{} \times 	\equivrep{g}{} &= \equivrep{(f \times g)}{}
\end{align*}
چې $(f + g)(n) = f(n) + g(n)$ او $(f \times g)(n) = f(n) \times
g(n)$۔ البته بايد دا هم وارزوو چې که $f$ او $g$ د کوشي لړۍ وي، نو
هره يوه $(f + g)$، $(f-g)$ او $(f\times g)$ هم د کوشي لړۍ ده؛ او
داسې ده، خو ثبوت ئې تاسو ته پرېږدو۔

په پاے کښې د مرتبې يو مفهوم تعريفوو۔ ووايئ $\equivrep{f}{}$
\emph{مثبت} دے هله او يوازې هله چې هم $\equivrep{f}{}\neq 0_\Real$
وي او هم $(\exists \ell \in \Nat)(\forall n > \ell)0 < f(n)$ وي۔
\textbf{د ژباړې د نسخې سمون (OLARI-012):} سرچينه په دې شرط کښې
ناطق صفر ليکي، خو کيڼه خوا د حقيقي عدد ټولګے دے او همدلته پورته
ثابته ناطقه تابع حقيقي عدد ته ورګډه شوې؛ نو د ډول له مخې سم صفر
حقيقي صفر دے۔ بيا ووايئ
$\equivrep{f}{} < \equivrep{g}{}$ هله او يوازې هله چې
$\equivrep{(g - f)}{}$ مثبت وي۔ بايد وارزوو چې دا سم تعريف شوے
دے (يعنې د هم ارزښتۍ له ټولګي څخه د «نماينده» تابعې پر ټاکنه پورې
اړه نۀ لري)۔
%\begin{align*}
%	\equivrep{f}{} \leq \equivrep{g}{} \text{ iff }&\equivrep{f}{} = \equivrep{g}{} \text{ or }\\
%	&(\exists \ell \in \Nat)(\forall n \in \Nat)(\ell < n \lif f(n) \leq g(n))
%\end{align*}
خو له دې ارزونې وروسته په ډېرې اسانۍ ښودل کېدے شي چې دا تعريفونه
سم الجبري خاصيتونه راکوي؛ يعنې:
\begin{thm}\ollabel{thm:cauchyorderedfield}
د کوشي د لړيو د هم ارزښتۍ ټولګي يو مرتب ميدان جوړوي۔
\end{thm}
\textbf{د ژباړې د نسخې سمون (OLARI-013):} انګرېزي سرچينه دلته او
د بشپړتيا په راتلونکې قضيه کښې په لنډه «د کوشي لړۍ» وايي؛ په لفظي
ډول يوازې خامې لړۍ ميدان نۀ جوړوي۔ دا پايلې د تعريف شوې اړيکې له مخې
د هغو د هم ارزښتۍ ټولګيو په اړه دي، لکه مخکښې تعريف شول۔

\begin{proof}
تمرين۔
\end{proof}

\begin{prob}
ثابته کړئ چې د کوشي د لړيو د هم ارزښتۍ ټولګي يو مرتب ميدان جوړوي۔
\end{prob}

دا ثابتول لږ ګران دي چې په دې ډول جوړ شوي حقيقي عددونه د بشپړتيا
خاصيت لري، نو ثبوت ئې وړاندې کوو۔

\begin{thm}
د کوشي د لړيو د هم ارزښتۍ د ټولګيو هر ناتش سټ چې پاسنی حد لري،
تر ټولو لږ پاسنی حد هم لري۔
\end{thm}

\begin{proof}[د ثبوت خاکه] فرض کړئ $S$ د کوشي د لړيو هر داسې ناتش
سټ دے چې پاسنی حد لري۔ نو داسې $p \in \Rat$ شته چې $p_{\Real}$ د
$S$ پاسنی حد دے۔ يو $r \in S$ واخلئ؛ بيا داسې $q \in \Rat$ شته
چې $q_{\Real} < r$ وي۔ نو که د $S$ تر ټولو لږ پاسنی حد موجود وي،
د $q_\Real$ او $p_\Real$ تر منځ (دواړه سرونه پکې شامل) دے۔
\textbf{د ژباړې د نسخې سمون (OLARI-014):} له يوه خپل‌سري حقيقي
پاسني حد څخه د ناطق پاسني حد، او د ټاکلي غړي څخه لاندې د ناطق
حد شتون د ناطقو عددونو د ارشيميدي/کثافت خاصيت ته اړتيا
لري۔ سرچينه دا منځنی دليل نۀ وړاندې کوي؛ لاندې نيمه‌کولو جوړونه
همدا ناطق سرونه فرضوي۔

موږ به تر ټولو لږ پاسني حد ته په يوه وخت له لاندې او پورته څخه
ورنژدې شو۔ په ځانګړې توګه دوه تابعې $f, g \colon \Nat \to \Rat$
داسې تعريفوو چې موخه مو دا وي چې $f$ له پورته او $g$ له لاندې څخه
تر ټولو لږ پاسني حد ته ورنژدې شي۔ په دې تعريف پيل کوو:
\begin{align*}
	f(0) &= p \\
	g(0) &= q
\end{align*}
بيا، چې $a_n = \frac{f(n) + g(n)}{2}$ وي، وټاکئ:\footnote{دا يو
بازګشتي تعريف دے۔ خو تر اوسه مو هېڅ دليل نۀ دے ورکړے چې
بازګشتي تعريفونه سم دي۔}
\begin{align*}
	f(n+1) &=
	\begin{cases}
		a_n &\text{که }(\forall h \in S)\equivrep{h}{} \leq (a_n)_\Real\\
		f(n)&\text{که نۀ}
	\end{cases}\\
	g(n+1) &=
	\begin{cases}
		a_n &\text{که }(\exists h \in S)\equivrep{h}{} \geq (a_n)_\Real\\
	 	g(n) &\text{که نۀ}
	\end{cases}
\end{align*}
$f$ او $g$ دواړه د کوشي لړۍ دي۔ (دا په نسبتاً اسانۍ ارزول کېدے
شي، خو د تمرين په توګه ئې پرېږدو۔) پام وکړئ چې تابع $(f-g)$ $0$ ته
ورنژدې کېږي، ځکه د $f$ او $g$ تر منځ توپير په هر ګام کښې نيم کېږي۔
نو $\equivrep{f}{} = \equivrep{g}{}$۔

لومړے ښيو چې $\equivrep{f}{}$ د $S$ پاسنی حد دے، يعنې
$(\forall h \in S)\equivrep{h}{} \leq \equivrep{f}{}$۔
(په لار کښې به \olref{thm:cauchyorderedfield} کاروو۔) يو $h \in S$
واخلئ او د نقيض دپاره فرض کړئ چې
$\equivrep{f}{} < \equivrep{h}{}$، نو
$0_\Real < \equivrep{(h-f)}{}$۔ ځکه $f$ يوه يکنواخته کموونکې د
کوشي لړۍ ده، داسې $n \in \Nat$ شته چې
$\equivrep{(c_{f(n)} - f)}{} < \equivrep{(h-f)}{}$۔ نو:
\[
	(f(n))_\Real = \equivrep{c_{f(n)}}{} < \equivrep{f}{} + \equivrep{(h-f)}{} = \equivrep{h}{},
\]
خو دا له دې جوړښتي حقيقت سره په ټکر کښې دے چې
$\equivrep{h}{} \leq (f(n))_\Real$۔

اوس ښيو چې $\equivrep{f}{} = \equivrep{g}{}$ د $S$ \emph{تر ټولو
لږ} پاسنی حد دے۔ نو $j$ دې هره د کوشي لړۍ وي او فرض کړئ
$\equivrep{j}{} < \equivrep{g}{}$۔ لکه پورته (د دې حقيقت په کارولو
چې $g$ \emph{زياتېدونکې} ده) داسې $n \in \Nat$ شته چې
$\equivrep{j}{} < (g(n))_\Real$۔ خو د جوړونې له مخې داسې $h \in S$
شته چې $(g(n))_\Real \leq \equivrep{h}{}$، نو
$\equivrep{j}{} < \equivrep{h}{}$ او له دې امله
$\equivrep{j}{}$ د $S$ پاسنی حد نۀ دے۔
\end{proof}

\end{document}
